% ============================================================
\chapter{Cơ sở lý thuyết}
\label{chap:background}
% ============================================================

\section{Truyền thông Backscatter}
\label{sec:bg:backscatter}

\subsection{Nguyên lý hoạt động của Backscatter}
\label{subsec:bg:backscatter:principle}
% Phân biệt passive tag và active reader.
% Hệ số phản xạ alpha, modulation bằng load impedance.

\subsection{Phân loại: Monostatic, Bistatic, Ambient}
\label{subsec:bg:backscatter:types}
% Mô tả và so sánh ba paradigm.
% Giải thích lý do chọn Ambient (AmBC) cho bài toán này.

\subsection{Mô hình tín hiệu Ambient Backscatter}
\label{subsec:bg:backscatter:signal}
% Phương trình tín hiệu nhận tại DU trong điều kiện jamming.
% Phân tích SINR cho backscatter.

\section{Tấn công Gây nhiễu vô tuyến}
\label{sec:bg:jamming}

\subsection{Phân loại kỹ thuật jamming}
\label{subsec:bg:jamming:types}
% Spot jammer, sweep jammer, barrage jammer, smart jammer.

\subsection{Mô hình jammer trong bài toán}
\label{subsec:bg:jamming:model}
% Ground-level constant-power jammer, ảnh hưởng đến SINR.

\subsection{Kỹ thuật chống gây nhiễu truyền thống}
\label{subsec:bg:antijamming:traditional}
% Frequency hopping, spread spectrum, power control.
% Hạn chế: phản ứng chậm, không tận dụng năng lượng jamming.

\section{Truyền thông D2D}
\label{sec:bg:d2d}

\subsection{Kiến trúc và ưu điểm D2D}
\label{subsec:bg:d2d:arch}

\subsection{Thách thức D2D dưới môi trường jamming}
\label{subsec:bg:d2d:challenge}

\section{UAV trong mạng truyền thông không dây}
\label{sec:bg:uav}

\subsection{Phân loại và đặc tính UAV}
\label{subsec:bg:uav:classification}

\subsection{Mô hình kênh Air-to-Ground}
\label{subsec:bg:uav:channel}
% Probabilistic LoS model (ITU-R), tổn hao trung bình.

\subsection{UAV làm relay động}
\label{subsec:bg:uav:relay}
% Decode-and-forward relay, energy model, kinematics.

\section{Học tăng cường}
\label{sec:bg:rl}

\subsection{Quá trình Quyết định Markov}
\label{subsec:bg:rl:mdp}
% State, action, reward, transition, discount factor.

\subsection{Deep Q-Network}
\label{subsec:bg:rl:dqn}

\subsection{Actor-Critic và DDPG}
\label{subsec:bg:rl:ddpg}
% Policy gradient, deterministic policy, replay buffer, target network.

\subsection{TD3}
\label{subsec:bg:rl:td3}
% Clipped double Q-learning, delayed policy update, target smoothing.

\section{Học tăng cường Đa tác tử}
\label{sec:bg:marl}

\subsection{Dec-POMDP và CTDE}
\label{subsec:bg:marl:decpomdp}

\subsection{MADDPG}
\label{subsec:bg:marl:maddpg}
% Centralized critic, decentralized actor.

\subsection{Imitation Learning và Behavior Cloning}
\label{subsec:bg:marl:il}
% Supervised BC loss, lambda annealing, expert warm-up.

\section{Transformer và Graph Attention Network}
\label{sec:bg:transformer-gat}

\subsection{Cơ chế Self-Attention và Transformer Encoder}
\label{subsec:bg:transformer}

\subsection{Graph Neural Network và Graph Attention}
\label{subsec:bg:gat}
% Multi-head attention trên đồ thị, message passing.

\subsection{Transformer-GAT trong học tăng cường đa tác tử}
\label{subsec:bg:tmac}
% TMAC-style encoder; lý do kết hợp Transformer + GAT cho multi-UAV.

\section{Tổng kết chương}
\label{sec:bg:summary}
